\begin{WidePar}
\begin{multicols}{3}
[\noindent\Ca{Die Auflistung dieser Lizenzen bedeutet \underline{nicht}, dass die
\AASS{} als Ganzes unter diesen Lizenzen lizensiert sind!} ]

\begin{description}
  \item[CC]
  \emph{Creative Commons}. Sammelbegriff für die Lizenzen der
  Creative Commons-Initiative, ohne nähere Angabe der Sublizenz.
  
  \item[CC-BY]
  \Speech{Creative-Commons Attribution Unported}
  \begin{inparadesc}
    \item[Lizenztext (Link):]
    Version 1.0: \url{http://creativecommons.org/licenses/by/1.0/};
    Version 2.0: \url{http://creativecommons.org/licenses/by/2.0/};
    Version 2.1: \url{http://creativecommons.org/licenses/by/2.1/};
    Version 2.5: \url{http://creativecommons.org/licenses/by/2.5/};
    Version 3.0: \url{http://creativecommons.org/licenses/by/3.0/}.
    \item[Kurzfassung (ohne Gewähr):] "`You let others copy, distribute,
    display, and perform your copyrighted work — and derivative works based
    upon it — but only if they give credit the way you request."'
  \end{inparadesc}
  
  \item[CC-BY-1.0]
  \Speech{Creative-Commons Attribution Unported},    % default = name
  \emph{"`Creative-Commons Attribution Unported 1.0."'}.
  \item[CC-BY-2.0]
  \Speech{Creative-Commons Attribution Unported},    % default = name
  \emph{"`Creative-Commons Attribution Unported 2.0."'}.
  \item[CC-BY-2.1]
  \Speech{Creative-Commons Attribution Unported},    % default = name
  \emph{"`Creative-Commons Attribution Unported 2.1."'}.
  \item[CC-BY-.at-3.0]
  \Speech{Creative-Commons Namensnennung 3.0 Österreich},
  \emph{Creative-Commons Namensnennung 3.0
  Österreich}. {\begin{inparadesc}
\item[Lizenztext (Link):]
\url{http://creativecommons.org/licenses/by/3.0/at/}
\item[Kurzfassung (ohne Gewähr):] \emph{Sie dürfen:}
\begin{inparaitem}
  \item das Werk vervielfältigen, verbreiten und öffentlich zugänglich machen
  \item Bearbeitungen des Werkes anfertigen;
\end{inparaitem}
\emph{Zu den folgenden Bedingungen:}
\begin{inparaitem}
  \item  Namensnennung: Sie müssen den Namen des Autors/Rechteinhabers in
  der von ihm festgelegten Weise nennen.
\end{inparaitem}
\end{inparadesc}}
\item[CC-BY-NC]
\Speech{Creative-Commons Attribution Non-Commercial},
\emph{Creative-Commons Attribution Non-Commercial}.
\begin{inparadesc}
  \item[Lizenztext (Link):]
  \url{http://creativecommons.org/licenses/by-nc/3.0/legalcode};
  \item[Kurzfassung (ohne Gewähr):] "`"'
\end{inparadesc}
\item[CC-BY-NC-ND]
\Speech{Creative-Commons Attribution Non-Commercial No Derivatives},
\emph{Creative-Commons Attribution Non-Commercial No Derivatives}.
\begin{inparadesc}
  \item[Lizenztext (Link):]
  \url{http://creativecommons.org/licenses/by-nc-nd/3.0/legalcode};
  \item[Kurzfassung (ohne Gewähr):] "`"'
\end{inparadesc}
\item[CC-BY-ND]
\Speech{Creative-Commons Attribution No Derivatives},
\emph{Creative-Commons Attribution No Derivatives}.
\begin{inparadesc}
  \item[Lizenztext (Link):]
  \url{http://creativecommons.org/licenses/by-nd/3.0/legalcode};
  \item[Kurzfassung (ohne Gewähr):] "`"'
\end{inparadesc}
\item[PD]
\emph{Public domain}, \emph{"`Gemeinfrei"'}.
Die Gemeinfreiheit bezeichnet alle Werke, die keinem Urheberrecht mehr unterliegen
oder ihm nie unterlegen haben. Das im angloamerikanischen Raum anzutreffende
Public Domain ist ähnlich, aber nicht identisch mit der europäischen
Gemeinfreiheit.
\end{description}
\end{multicols}
\end{WidePar}